\section{Introduction}

Welcome and thank you for your interest in the \Grackle{} (\grackle{}) at Trinity University!
This document summarizes the overarching goals and philosophy of the laboratory, recommendations on how to get involved, my expectations for lab members, and some reasonable expectations lab members should have for me.

As the Primary Investigator (PI) of the lab, I\footnote{\theauthor} wrote this document, but use the pronoun ``we'' throughout to reflect the collaborative nature of our work and my intentions to support collective work toward shared goals.

\subsection{Research goals, approaches, \& philosophy}

As the laboratory name implies, we use \emph{geochemistry}---with some special emphasis on \emph{radioisotopes}---to reconstruct both time and process in Earth's history.
Researchers in the  group use a variety of techniques to answer interdisciplinary questions about the recent history of ice and water on/near our planet's surface and the earliest stages of planetary accretion and evolution in our solar system.  
We use both \emph{analytical} chemistry techniques---to measure elemental and isotopic compositions of Earth and planetary material---as well as \emph{computational} techniques---to develop numerical models that help us interpret Earth and planetary processes from chemical data.
In plainer language, we measure the chemistry of Earth/planetary materials and write computer codes to simulate the natural systems that resulted in the chemical compositions we/others have measured. 


Members of this lab rely on principles of (geo)chemistry to interpret Earth and planetary histories.
In many cases, we use analytical and wet chemistry methods to generate novel data. 
\emph{In all cases}, we write and implement scientific computer code to visualize data and model natural systems.
Yet, we also embrace the aphorism\footnote{Commonly attributed to statistician George E.P.~Box. However, folks have been saying various forms of this statement for a very long time.} 
``All models are wrong, but some are useful.''
In other words, we strive to consider the limitations of our models, recognizing not only what they \emph{do} tell us, but what they \emph{do not} tell us.

My research interests and goals are broad, and I encourage you to lean into the topics that you find interesting within the Earth and planetary sciences. 
If you have an idea or a topic you are interested in exploring, let's talk about it! 
I cannot guarantee that I will be able or willing to support your research on that topic, but I will fully consider the potential.